\documentclass[12pt]{article}

% setting margins for the document
\usepackage[letterpaper, margin=1.1in]{geometry}

\usepackage[T1]{fontenc}
% packages used for the various math commands
\usepackage{amsmath, amsthm, amssymb,enumitem}

% setting up the font
\usepackage{libertinus-type1}
\usepackage{libertinust1math}

% adjusting in-word spacing
\usepackage{microtype}

\usepackage{parskip}
\usepackage{etoolbox}

% setting custom margins for enumerate and itemizaed lists
\setlist[enumerate]{parsep=0pt plus 4pt,topsep=3pt plus 4pt}
\setlist[itemize]{parsep=0pt plus 4pt,topsep=3pt plus 4pt,itemsep=.5ex}

\leftmargini=5.5ex
\leftmarginii=3.5ex


%For separated lists with consecutive numbering
\newcounter{separated}
\newcommand{\excise}[1]{}
\newcommand{\comment}[1]{{$\star$\sf\textbf{#1}$\star$}}
\newcommand{\exnumber}[1]{\noindent\textbf{#1}}

%new math symbols taking no arguments
\newcommand\0{\mathbf{0}}
\newcommand\CC{\mathbb{C}}
\newcommand\FF{\mathbb{F}}
\newcommand\NN{\mathbb{N}}
\newcommand\QQ{\mathbb{Q}}
\newcommand\RR{\mathbb{R}}
\newcommand\ZZ{\mathbb{Z}}
\newcommand\bb{\mathbf{b}}
\newcommand\kk{\Bbbk}
\newcommand\mm{\mathfrak{m}}
\newcommand\xx{\mathbf{x}}
\newcommand\yy{\mathbf{y}}
\newcommand\GL{\mathit{GL}}
\newcommand\pset{\mathcal{P}}
\newcommand\into{\hookrightarrow}
\newcommand\nsub{\trianglelefteq}
\newcommand\onto{\twoheadrightarrow}
\newcommand\minus{\smallsetminus}
\newcommand\goesto{\rightsquigarrow}

% theorem environment
\newtheorem*{theorem}{Theorem}

%redefined math symbols taking no arguments
\newcommand\<{\langle}
\renewcommand\>{\rangle}
\renewcommand\iff{\Leftrightarrow}
\renewcommand\phi{\varphi}
\renewcommand\implies{\Rightarrow}

%new math symbols taking arguments
\newcommand\ol[1]{{\overline{#1}}}

%redefined math symbols taking arguments
\renewcommand\mod[1]{\ (\mathrm{mod}\ #1)}

%roman font math operators
\DeclareMathOperator\im{im}

%for easy 2 x 2 matrices
\newcommand\twobytwo[1]{\left[\begin{array}{@{}cc@{}}#1\end{array}\right]}

%for easy column vectors of size 2
\newcommand\tworow[1]{\left[\begin{array}{@{}c@{}}#1\end{array}\right]}

% redefining \maketitle so that it has proper spacing
\makeatletter
\renewcommand{\maketitle}{%
  \begingroup
  \setlength{\parskip}{0pt}   % kills the parskip effect inside the title block
  \begin{center}
    {\LARGE\bfseries \@title\par}
    \vspace{1.5ex}
    {\large \@author\par}
    \vspace{0.75ex}
    {\normalsize \@date\par}
  \end{center}
  \endgroup
  \vspace{2ex}
}
\makeatother

% makes sure that proofs within an itemized or enumerate environment 
% maintain proper spacing between problem statements. 
\makeatletter
\BeforeBeginEnvironment{proof}{%
  \par
  \ifnum\@listdepth>0 \addvspace{\medskipamount}\fi
}
\AfterEndEnvironment{proof}{%
  \ifnum\@listdepth>0 \addvspace{\medskipamount}\fi
}
\makeatother

%%%%%%%%%%%%%%%%%%%%%%%%%%%%%%%%%%%%%%%%%%%%%%%%%%%%%%%%%%%%%%%%%%%%%%
\begin{document}%%%%%%%%%%%%%%%%%%%%%%%%%%%%%%%%%%%%%%%%%%%%%%%%%%%%%%
%%%%%%%%%%%%%%%%%%%%%%%%%%%%%%%%%%%%%%%%%%%%%%%%%%%%%%%%%%%%%%%%%%%%%%
	\title{Countable and Uncountable Sets Solutions}
	\author{Solutions by: Raindrops\_drop}
	\date{October 11, 2026}
	\maketitle
	
	\section*{Exercises}

	\exnumber{1.} Show that $\QQ$ is countably infinite.

	\begin{proof}
		Recall $\ZZ$ is countably infinite. Note $\ZZ - \{0\}$
		is an infinite subset of $\ZZ$, and thus by Theorem 7.3 the set 
		$\ZZ - \{0\}$ is countable.
		Furthermore, the set $\ZZ \times (\ZZ - \{0\})$ is also countable
		by Theorem 7.6. Thus, define $f: \ZZ \times (\ZZ - \{0\}) \to \QQ$
		by the equation $f(n, m) = n/m$. Note that $f$ is a surjection
		since for any $a/b \in \QQ$, $a \in \ZZ$ and $b \in \ZZ - \{0\}$.
		Given that $\ZZ \times (\ZZ - \{0\})$ is countable, there exists
		a surjection $g: \ZZ_{+} \to \ZZ \times (\ZZ - \{0\})$.
		As such, $f \circ g: \ZZ_{+} \to \QQ$ is a surjection. Therefore,
		$\QQ$ is countable. Since $\ZZ \subset \QQ$, we further conclude
		$\QQ$ is countably infinite. 
	\end{proof}

	\exnumber{2.} Show that the maps $f$ and $g$ of Examples 1 and 2
	are bijections. 

	\begin{proof}
		The map $f: \ZZ \to \ZZ_{+}$ in Example 1 is defined as 
		$$
		f(n) = 
		\begin{cases}
			2n & n > 0, \\
			-2n + 1 & n \leq 0.
		\end{cases}
		$$
		To show $f$ is an injective map, let $m, n \in \ZZ$ such
		that $f(n) = f(m) = a$. In particular, note $m$ and 
		$n$ must either be both positive or both non-positive. This is 
		because if without loss of generality $n$ is positive and $m$
		is non-positive we have $f(n) = 2n$ and $f(m) = -2m + 1$. 
		In particular, $2n = -2m + 1 = a$, which cannot be true
		since $2n$ is even and $-2m + 1$ is odd, and a positive integer
		cannot be both even and odd. Thus, if both $n$ and $m$ are positive,
		then $2n = 2m$ by assumption, which implies $n = m$.
		Hence, $f$ is injective. If $n, m$ are both non-positive,
		then $-2n + 1 = -2m + 1$. This implies $-2n = -2m$, and 
		thus $n = m$, which allows us to conclude $f$ is injective.

		To show $f$ is surjective, consider $a \in \ZZ_{+}$. If 
		$a$ is even, $a = 2n$ for some $n \in \ZZ_{+}$, 
		as such there exists such $n \in \ZZ$ that is positive and is 
		mapped by $f$ onto $a$. If $a$ is odd, let $n = (1 - a)/2$.
		Since $a$ is odd, $1 - a$ is equivalent to $1 - 2k - 1$ for 
		some $k \in \ZZ$, and thus $1 - a$ is an even integer,
		implying $(1 - a)/2$ is also an integer. Since $a \geq 1$,
		$1 - a \leq 0$, implying $(1-a)/2 \leq 0$ also. Thus,
		$n \in \ZZ$ is a non-positive integer which is mapped by $f$
		onto $a$. Therefore, $f$ is surjective, and thus we conclude 
		$f$ is a bijection.

		The map $g: A \to \ZZ_{+}$ is defined 
		by the formula
		$$
		g(n, m) = \frac{1}{2}(n - m)n + m,
		$$
		where $A$ is the  subset of $\ZZ_{+} \times \ZZ_{+}$ consisting
		of pairs $(n, m)$ for which $m \leq n$.  
		\comment{Haven't proved that this g function is a bijection.}
	\end{proof}

	\exnumber{3.} Let $X$ ve the two-element set $\{0, 1\}$. Show there is 
	a bijective correspondence between the set $\pset(\ZZ_{+})$ abd the 
	cartesian product $X^\omega$.

	\exnumber{4.}
	\begin{enumerate}[label=(\alph*)]
		\item A real number $x$ is said to be \textbf{algebraic} 
		(over the rationals) if it satisfies some polynomial equation 
		of positive degree 
		$$
		x^n + a_{n - 1}x^{n - 1} + \cdots + a_1 x + a_0 = 0
		$$
		with rational coefficients $a_i$. Assuming that each polynomial
		equation has only finitely many roots, show that the set 
		of algebraic numbers is countable. 
		\item A real number is said to be \textbf{transcendental} if it is 
		not algebraic. Assuming that the reals are uncountable, show that 
		the transcendental numbers are uncountable. (It is a somewhat 
		surprising fact that only two transcendental numbers are familiar
		to us: $e$ and $\pi$. Even proving these two numbers transcendental
		is highly non-trivial).
	\end{enumerate}

	\exnumber{5.} Determine, for each of the following sets, whether or 
	not it is countable. Justify your answers.
	\begin{enumerate}[label=(\alph*)]
		\item The set $A$ of all functions $f: \{0, 1\} \to \ZZ_{+}$.
		\item The set $B_n$ of all functions $f: \{1, \ldots, n\} \to \ZZ_{+}$.
		\item The set $C = \bigcup_{n \in \ZZ_{+}} B_n$.
		\item The set $D$ of all functions $f: \ZZ_{+} \to \ZZ_{+}$.
		\item The set $E$ of all functions $f: \ZZ_{+} \to \{0, 1\}$.
		\item The set $F$ of all functions $f: \ZZ_{+} \to \{0, 1\}$ that
		are "eventually zero." [We say that $f$ is \textbf{eventually zero} 
		uf there is a positive integer $N$ such that $f(n) = 0$ for all
		$n \geq N$].
		\item The set $G$ of all functions $f: \ZZ_{+} \to \ZZ_{+}$ that 
		are eventually 1.
		\item The set $H$ of all functions $f: \ZZ_{+} \to \ZZ_{+}$ that
		are eventually constant. 
		\item The set $I$ of all two-element subsets of $\ZZ_{+}$.
		\item The set $J$ of all finite subsets of $\ZZ_{+}$. 
	\end{enumerate}

	\exnumber{6.} We say that two sets $A$ and $B$ \textbf{have the same
	cardinality} if there is a bijection of $A$ with $B$.
	\begin{enumerate}[label=(\alph*)]
		\item Show that if $B \subseteq A$ and if there is an injection
		$$
		f: A \into B,
		$$
		then $A$ and $B$ have the same cardinality. [Hint: Define $A_1 = A$,
		$B_1 = B$, and for $n > 1$, $A_n = f(A_{n - 1})$ and $B_n = f(B_{n - 1})$.
		Note that $A_1 \supseteq B_1 \supseteq A_2 \supseteq B_2 \supseteq \cdots$.
		Define a bijection $h: A \to B$ by the rule
		$$
		h(x) = 
		\begin{cases}
			f(x) & \text{if } x \in A_n - B_n \text{ for some } n, \\
			x & \text{otherwise.}]
		\end{cases}
		$$
		\item 
		\begin{theorem}
			(Schroeder-Bernstein Theorem) If there are injections $f: 
			A \into C$ and $g: C \into A$, then $A$ and $C$ have the 
			same cardinality.  	
		\end{theorem}
	\end{enumerate}

	\exnumber{7.} Show that the sets $D$ and $E$ of Exercise 5. have the 
	same cardinality.

	\exnumber{8.} Let $X$ denote the two-element set $\{0, 1\}$;
	let $\mathcal{B}$ be the set of countable subsets of $X^\omega$.
	Show that $X^\omega$ and $\mathcal{B}$ have the same cardinality.

	\exnumber{9.} 
	\begin{enumerate}[label=(\alph*)]
		\item The formula
		$$
		\begin{aligned}
			h(1) &= 1  \\
			h(2) &= 2  \\
			h(n) &= [h(n + 1)]^2 - [h(n - 1)]^2 \hspace{5mm} n \geq 2
		\end{aligned}
		$$
		is not one to which the principle of recursive definition 
		applies. Show that nevertheless there does exists a function 
		$h: \ZZ_{+} \to \RR$ satisfying this formula. [Hint: Reformulate
		the formula so that the principle will apply and require $h$ to 
		be positive].

		\item Show that the formula in part (a) does not determine $h$
		uniquely. [Hint: If $h$ is a positive function satisfying the formula 
		let $f(i) = h(i)$ for $i \neq 3$, and let $f(3) = -h(3)$].
		
		\item Show that there is no function $f: \ZZ_{+} \to \RR$ satisfying
		the formula 
		$$
		\begin{aligned}
			h(1) &= 1  \\
			h(2) &= 2  \\
			h(n) &= [h(n + 1)]^2 + [h(n - 1)]^2 \hspace{5mm} n \geq 2
		\end{aligned}
		$$
	\end{enumerate}
%%%%%%%%%%%%%%%%%%%%%%%%%%%%%%%%%%%%%%%%%%%%%%%%%%%%%%%%%%%%%%%%%%%%%%
\end{document}%%%%%%%%%%%%%%%%%%%%%%%%%%%%%%%%%%%%%%%%%%%%%%%%%%%%%%%%
%%%%%%%%%%%%%%%%%%%%%%%%%%%%%%%%%%%%%%%%%%%%%%%%%%%%%%%%%%%%%%%%%%%%%%
